\documentclass[12pt,a4paper]{article}

% ── Encodage (pdflatex) ─────────────────────────────────────────────────────
\usepackage[utf8]{inputenc}
\usepackage[T1]{fontenc}
\usepackage{lmodern}
\usepackage{amssymb}

% ── Mise en page ────────────────────────────────────────────────────────────
\usepackage[top=2.5cm, bottom=2.5cm, left=2.8cm, right=2.8cm]{geometry}
\usepackage{parskip}
\usepackage{microtype}
\usepackage{setspace}
\onehalfspacing

% ── Couleurs et boites ──────────────────────────────────────────────────────
\usepackage[dvipsnames]{xcolor}
\usepackage{tcolorbox}
\tcbuselibrary{skins, breakable, listings}

\definecolor{codebg}{HTML}{1E1E2E}
\definecolor{codefg}{HTML}{CDD6F4}
\definecolor{commentfg}{HTML}{7F849C}
\definecolor{keywordfg}{HTML}{CBA6F7}
\definecolor{stringfg}{HTML}{A6E3A1}
\definecolor{numberfg}{HTML}{FAB387}
\definecolor{framebg}{HTML}{F5F5F5}
\definecolor{frameaccent}{HTML}{204A87}
\definecolor{alertbg}{HTML}{FFF3CD}
\definecolor{alertborder}{HTML}{D6940B}
\definecolor{okbg}{HTML}{D4EDDA}
\definecolor{okborder}{HTML}{28A745}
\definecolor{errorbg}{HTML}{F8D7DA}
\definecolor{errorborder}{HTML}{DC3545}

% ── Listings ────────────────────────────────────────────────────────────────
\usepackage{listings}

% Caracteres accentues dans les commentaires des blocs de code
\lstset{literate=
  {a}{{a}}1 {e}{{e}}1 {i}{{i}}1 {o}{{o}}1 {u}{{u}}1
  {A}{{A}}1 {E}{{E}}1 {I}{{I}}1 {O}{{O}}1 {U}{{U}}1
  {\'e}{{\'{e}}}1 {\`e}{{\`{e}}}1 {\^e}{{\^{e}}}1 {\"e}{{\"e}}1
  {\'a}{{\'{a}}}1 {\`a}{{\`{a}}}1 {\^a}{{\^{a}}}1
  {\`u}{{\`{u}}}1 {\^u}{{\^{u}}}1 {\"u}{{\"u}}1
  {\^i}{{\^{i}}}1 {\"i}{{\"i}}1
  {\^o}{{\^{o}}}1 {\"o}{{\"o}}1
  {\c c}{{\c{c}}}1 {\c C}{{\c{C}}}1
  {->}{{$\rightarrow$}}2
  {>=}{{$\geq$}}2
  {<=}{{$\leq$}}2
}

\lstdefinestyle{cppstyle}{
  backgroundcolor=\color{codebg},
  basicstyle=\ttfamily\small\color{codefg},
  commentstyle=\color{commentfg}\itshape,
  keywordstyle=\color{keywordfg}\bfseries,
  stringstyle=\color{stringfg},
  numberstyle=\tiny\color{numberfg},
  breakatwhitespace=false,
  breaklines=true,
  keepspaces=true,
  numbers=left,
  numbersep=10pt,
  showspaces=false,
  showstringspaces=false,
  tabsize=4,
  frame=none,
  xleftmargin=15pt,
  language=C++,
  morekeywords={uint8_t,uint16_t,uint32_t,Serial1,delay,memcpy,atoi},
  inputencoding=utf8,
  extendedchars=true,
}

\newtcblisting{cppbox}[1][]{
  colback=codebg,
  colframe=frameaccent,
  arc=4pt,
  boxrule=0.8pt,
  left=6pt, right=6pt, top=4pt, bottom=4pt,
  listing only,
  listing style=cppstyle,
  breakable,
  #1
}

\newtcolorbox{infobox}[1][]{
  colback=framebg, colframe=frameaccent,
  arc=4pt, boxrule=0.8pt,
  fonttitle=\bfseries, #1
}
\newtcolorbox{alertbox}[1][]{
  colback=alertbg, colframe=alertborder,
  arc=4pt, boxrule=0.8pt,
  fonttitle=\bfseries, #1
}
\newtcolorbox{okbox}[1][]{
  colback=okbg, colframe=okborder,
  arc=4pt, boxrule=0.8pt,
  fonttitle=\bfseries, #1
}
\newtcolorbox{errorbox}[1][]{
  colback=errorbg, colframe=errorborder,
  arc=4pt, boxrule=0.8pt,
  fonttitle=\bfseries, #1
}

% ── Tableaux ────────────────────────────────────────────────────────────────
\usepackage{booktabs}
\usepackage{array}
\usepackage{longtable}
\newcolumntype{L}[1]{>{\raggedright\arraybackslash}p{#1}}
\newcolumntype{C}[1]{>{\centering\arraybackslash}p{#1}}

% ── En-tetes et pied de page ────────────────────────────────────────────────
\usepackage{fancyhdr}
\pagestyle{fancy}
\fancyhf{}
\fancyhead[L]{\small\textcolor{gray}{Drone \`a voile --- Module GPS}}
\fancyhead[R]{\small\textcolor{gray}{Documentation technique}}
\fancyfoot[C]{\small\thepage}
\renewcommand{\headrulewidth}{0.4pt}

\usepackage[colorlinks=true,linkcolor=frameaccent,urlcolor=frameaccent,
            pdftitle={Module GPS u-blox NEO-6M},
            pdfauthor={Facundo Arito}]{hyperref}

% ── Titres de section ────────────────────────────────────────────────────────
\usepackage{titlesec}
\titleformat{\section}{\Large\bfseries\color{frameaccent}}{}{0em}{\thesection\enspace}[\titlerule]
\titleformat{\subsection}{\large\bfseries\color{frameaccent!80}}{}{0em}{\thesubsection\enspace}
\titleformat{\subsubsection}{\normalsize\bfseries}{}{0em}{\thesubsubsection\enspace}

% ── Commandes utilitaires ────────────────────────────────────────────────────
\newcommand{\hex}[1]{\texttt{0x#1}}
\newcommand{\code}[1]{\texttt{#1}}

% ════════════════════════════════════════════════════════════════════════════
\begin{document}
% ════════════════════════════════════════════════════════════════════════════

% ── Page de titre ───────────────────────────────────────────────────────────
\begin{titlepage}
  \centering
  \vspace*{2cm}

  {\Huge\bfseries\color{frameaccent} Module GPS u-blox NEO-6M}\\[0.6em]
  {\Large Fonctionnement, diagnostic et correction}\\[2em]

  \textcolor{gray}{\rule{0.6\textwidth}{0.5pt}}\\[1.5em]

  {\large Projet \textbf{Drone \`a voile}\\[0.4em]
  LilyGO T-Beam V1.1 --- ESP32 --- Antenne Taoglas ADFGP.25A}\\[3em]

  \begin{tabular}{rl}
    \textbf{Auteur}         & Facundo Arito \\[0.3em]
    \textbf{Date}           & 20 mai 2026 \\[0.3em]
    \textbf{Fichier source} & \code{main/src/drivers/GpsUart.cpp} \\
  \end{tabular}

  \vfill
  \textcolor{gray}{\small Document g\'en\'er\'e automatiquement --- usage interne}
\end{titlepage}

\tableofcontents
\newpage

% ════════════════════════════════════════════════════════════════════════════
\section{Vue d'ensemble du syst\`eme GPS}
% ════════════════════════════════════════════════════════════════════════════

\subsection{R\^ole du GPS dans le drone}

Le module GPS fournit au drone sa position g\'eographique (latitude, longitude),
sa vitesse sur l'eau et son cap. Ces donn\'ees sont essentielles pour le pilotage
automatique et pour le journal de navigation.

\subsection{Composants mat\'eriels}

\begin{infobox}[title=Mat\'eriel impliqu\'e]
\begin{tabular}{L{4.2cm} L{9.5cm}}
\toprule
\textbf{Composant} & \textbf{R\^ole} \\
\midrule
LilyGO T-Beam V1.1
  & Carte principale --- int\`egre un ESP32, un module u-blox NEO-6M
    et un connecteur u.FL pour antenne externe \\[0.4em]
u-blox NEO-6M
  & R\'ecepteur GPS. Communique avec l'ESP32 via UART (\code{Serial1})
    et accepte des commandes de configuration binaires (protocole UBX) \\[0.4em]
Taoglas ADFGP.25A
  & Antenne active externe (patch c\'eramique + LNA int\'egr\'e).
    Le LNA n\'ecessite une tension DC transmise sur le c\^able coaxial \\[0.4em]
C\^able u.FL vers SMA
  & Connecte le NEO-6M \`a l'antenne externe \\
\bottomrule
\end{tabular}
\end{infobox}

\subsection{Flux de donn\'ees}

\begin{center}
\begin{tabular}{ccccc}
\textbf{Antenne} & $\longrightarrow$ & \textbf{NEO-6M} & $\longrightarrow$ & \textbf{ESP32} \\
Signal RF 1575\,MHz & & D\'ecode les satellites & & Phrases NMEA (ASCII) \\
\end{tabular}
\end{center}

\medskip
Le NEO-6M re\c{c}oit les signaux radio des satellites, calcule la position
et \'emet des \textbf{phrases NMEA} sur son UART. L'ESP32 lit ce flux via
\code{Serial1} et utilise la biblioth\`eque \code{TinyGPSPlus} pour en
extraire les donn\'ees utiles.

% ════════════════════════════════════════════════════════════════════════════
\section{Le protocole UBX}
% ════════════════════════════════════════════════════════════════════════════

\subsection{Deux langages sur un m\^eme fil}

Le NEO-6M supporte deux protocoles sur son UART~:

\begin{itemize}
  \item \textbf{NMEA} --- phrases textuelles (ex.~\code{\$GPGGA,...}) lues
        pour extraire la position. C'est le flux normal en fonctionnement.
  \item \textbf{UBX} --- protocole binaire propri\'etaire u-blox, utilis\'e pour
        \emph{configurer} le module (activer des fonctions, modifier des
        param\`etres, sauvegarder en m\'emoire).
\end{itemize}

Le NEO-6M distingue une trame UBX d'une phrase NMEA gr\^ace aux deux premiers
octets~: \hex{B5 62} pour UBX, le caract\`ere \code{\$} pour NMEA.

\subsection{Structure d'une trame UBX}

\begin{center}
\begin{tabular}{C{1.8cm} C{1.8cm} C{2cm} C{2cm} C{2.5cm} C{2.5cm}}
\toprule
\textbf{Sync 1} & \textbf{Sync 2} & \textbf{Classe} & \textbf{ID}
  & \textbf{Longueur (LE)} & \textbf{Payload} \\
\midrule
\hex{B5} & \hex{62} & 1 octet & 1 octet & 2 octets & $N$ octets \\
\bottomrule
\end{tabular}
\end{center}

\medskip
Chaque trame se termine par un \textbf{checksum de Fletcher} sur 2 octets
(\code{CK\_A}, \code{CK\_B}), calcul\'e sur tous les octets depuis la classe
jusqu'\`a la fin du payload (les deux octets de synchronisation sont exclus).

\subsection{Calcul du checksum de Fletcher}

\begin{cppbox}
uint8_t CK_A = 0, CK_B = 0;
// Parcourir : classe + ID + longueur (2 octets) + payload
for (int i = 0; i < N; i++) {
    CK_A += data[i];
    CK_B += CK_A;
}
// CK_A et CK_B sont appended a la fin de la trame
\end{cppbox}

Ce checksum garantit l'int\'egrit\'e de la trame~: si un bit est corrompu,
le NEO-6M ignore la commande.

% ════════════════════════════════════════════════════════════════════════════
\section{Pourquoi le GPS ne fonctionnait pas}
% ════════════════════════════════════════════════════════════════════════════

\subsection{Sympt\^ome observ\'e}

Le compteur \code{chars} de \code{TinyGPSPlus} restait \textbf{fig\'e \`a 10}
malgr\'e le GPS aliment\'e et l'antenne connect\'ee.

\begin{alertbox}[title={Diagnostic~: que signifie chars = 10~?}]
\code{chars} comptabilise \textbf{chaque octet re\c{c}u} sur \code{Serial1}.
Une trame UBX ACK-ACK (r\'eponse d'acquittement du NEO-6M) fait exactement
\textbf{10 octets}~:

\medskip
\texttt{B5 62 05 01 02 00 06 09 17 6B}

\medskip
\begin{tabular}{ll}
\texttt{B5 62} & Sync UBX \\
\texttt{05 01} & Classe ACK, ID ACK \\
\texttt{02 00} & Longueur payload = 2 \\
\texttt{06 09} & \'Echo de la commande acquitt\'ee (classe CFG, ID CFG) \\
\texttt{17 6B} & Checksum \\
\end{tabular}

\medskip
Conclusion~: le GPS \'etait \textbf{bien aliment\'e et fonctionnel},
il r\'epondait aux commandes UBX, mais n'\'emettait \textbf{aucune phrase NMEA}.
\end{alertbox}

\subsection{Cause racine~: la m\'emoire flash interne du NEO-6M}

Le NEO-6M poss\`ede une \textbf{m\'emoire flash interne} qui persiste la
configuration entre les mises sous tension. Une session de configuration
ant\'erieure (via le logiciel u-center de u-blox) avait sauvegard\'e une
configuration qui \textbf{d\'esactivait toutes les phrases NMEA}.

\begin{errorbox}[title=Cause identifi\'ee]
La configuration persistante en flash d\'esactivait la sortie NMEA sur le
port UART du NEO-6M. Le module fonctionnait normalement mais n'\'emettait que
des r\'eponses UBX, invisibles pour \code{TinyGPSPlus}.
\end{errorbox}

\subsection{Cause aggravante~: antenne active sans alimentation DC}

L'antenne Taoglas ADFGP.25A est une antenne \textbf{active}~: elle contient
un amplificateur LNA qui n\'ecessite une tension DC transmise sur le c\^able
coaxial. Par d\'efaut apr\`es une r\'einitialisation ROM, le superviseur
d'antenne du NEO-6M est \textbf{d\'esactiv\'e}~:

\begin{itemize}
  \item Aucune tension DC sur le connecteur u.FL $\Rightarrow$ LNA hors tension.
  \item Le commutateur RF du T-Beam pointe vers le patch c\'eramique interne
        au lieu de l'antenne externe.
  \item R\'esultat~: z\'ero signal GPS m\^eme avec l'antenne externe branch\'ee.
\end{itemize}

% ════════════════════════════════════════════════════════════════════════════
\section{Les deux commandes UBX de correction}
% ════════════════════════════════════════════════════════════════════════════

\subsection{Commande 1 --- CFG-CFG~: r\'einitialisation usine}

\subsubsection{Objectif}

Effacer la configuration sauvegard\'ee en flash et recharger les param\`etres
d'usine depuis la ROM interne. Apr\`es cette commande~:

\begin{itemize}
  \item Le UART1 du NEO-6M repasse \`a 9\,600 bauds.
  \item Toutes les phrases NMEA standard sont r\'eactiv\'ees.
  \item Le superviseur d'antenne est d\'esactiv\'e (\'etat ROM par d\'efaut).
\end{itemize}

\subsubsection{Structure de la trame}

\begin{center}
\begin{tabular}{C{1.8cm} C{1.5cm} C{1.5cm} C{1.8cm} C{2cm} C{1.5cm}}
\toprule
\textbf{Sync} & \textbf{Classe} & \textbf{ID} & \textbf{Longueur}
  & \textbf{Payload} & \textbf{CRC} \\
\midrule
\hex{B5 62} & \hex{06} & \hex{09} & \hex{0D 00} & 13 octets & \hex{71 FE} \\
\bottomrule
\end{tabular}
\end{center}

Classe \hex{06} = CFG (configuration), ID \hex{09} = sous-commande de gestion
de la configuration persistante.

\subsubsection{D\'ecomposition du payload (13 octets, little-endian)}

\begin{center}
\begin{tabular}{C{3.5cm} C{3cm} L{6.5cm}}
\toprule
\textbf{Champ (32\,bits LE)} & \textbf{Valeur hex} & \textbf{Signification} \\
\midrule
\code{clearMask} & \hex{0000001F} & Efface 5 sections de configuration \\
\code{saveMask}  & \hex{00000000} & Ne sauvegarde rien \\
\code{loadMask}  & \hex{0000001F} & Recharge depuis la ROM \\
\code{deviceMask} & \hex{17} (8\,bits) & Cible~: BBR + Flash + EEPROM \\
\bottomrule
\end{tabular}
\end{center}

\begin{infobox}[title=Pourquoi clearMask = 0x1F~?]
Ce champ est un masque de bits~: chaque bit correspond \`a une section
de la configuration persistante.

\medskip
\begin{tabular}{C{1.5cm} C{2.5cm} L{8cm}}
\toprule
\textbf{Bit} & \textbf{Valeur hex} & \textbf{Section effac\'ee} \\
\midrule
0 & \hex{01} & \code{ioPort}  --- configuration des ports (UART, SPI, USB) \\
1 & \hex{02} & \code{msgConf} --- activation des phrases NMEA/UBX par port \\
2 & \hex{04} & \code{infMsg}  --- messages d'information \\
3 & \hex{08} & \code{navConf} --- param\`etres de navigation (mode, SBAS) \\
4 & \hex{10} & \code{rxmConf} --- mode r\'ecepteur \\
\midrule
\multicolumn{3}{c}{\hex{01}\,|\,\hex{02}\,|\,\hex{04}\,|\,\hex{08}\,|\,\hex{10}
  = \hex{1F}} \\
\bottomrule
\end{tabular}

\medskip
La section \code{msgConf} (bit 1) est la plus critique~: c'est elle qui
d\'esactivait les phrases NMEA dans notre cas.
\end{infobox}

\subsubsection{Code source et explication octet par octet}

\begin{cppbox}
static const uint8_t kCfgCfgReset[] = {
    0xB5, 0x62,              // Sync UBX -- debut obligatoire de toute trame
    0x06, 0x09,              // Classe CFG (0x06), ID CFG (0x09)
    0x0D, 0x00,              // Longueur payload : 13 octets (little-endian)
    0x1F, 0x00, 0x00, 0x00, // clearMask : efface 5 sections de config
    0x00, 0x00, 0x00, 0x00, // saveMask  : ne sauvegarde rien
    0x1F, 0x00, 0x00, 0x00, // loadMask  : recharge depuis la ROM
    0x17,                    // deviceMask : BBR(1)+Flash(2)+EEPROM(4)+SPI(16)
    0x71, 0xFE               // Checksum Fletcher (CK_A=0x71, CK_B=0xFE)
};
\end{cppbox}

\begin{alertbox}[title=Pourquoi \texttt{static const}~?]
\code{static} dans une fonction signifie que le tableau est allou\'e \textbf{une
seule fois en m\'emoire flash} (section \code{.rodata}), et non sur la pile
\`a chaque appel. Sur un microcontr\^oleur avec peu de RAM comme l'ESP32,
c'est une bonne pratique pour les constantes volumineuses.
\end{alertbox}

\subsection{Commande 2 --- CFG-ANT~: activation de l'antenne active}

\subsubsection{Objectif}

Activer le \textbf{superviseur d'antenne} du NEO-6M, qui accomplit trois
actions~:

\begin{enumerate}
  \item Fournit une tension DC sur le c\^able coaxial pour alimenter le LNA
        de l'antenne Taoglas.
  \item Commande le commutateur RF du T-Beam (via la broche \code{ANT\_FLAG}
        du NEO-6M) pour s\'electionner l'antenne externe u.FL.
  \item Active la d\'etection de court-circuit pour prot\'eger le module.
\end{enumerate}

\subsubsection{Structure de la trame}

\begin{center}
\begin{tabular}{C{1.8cm} C{1.5cm} C{1.5cm} C{1.8cm} C{2cm} C{1.5cm}}
\toprule
\textbf{Sync} & \textbf{Classe} & \textbf{ID} & \textbf{Longueur}
  & \textbf{Payload} & \textbf{CRC} \\
\midrule
\hex{B5 62} & \hex{06} & \hex{13} & \hex{04 00} & 4 octets & \hex{47 57} \\
\bottomrule
\end{tabular}
\end{center}

Classe \hex{06} = CFG, ID \hex{13} = ANT (configuration de l'antenne).

\subsubsection{D\'ecomposition du payload}

\begin{center}
\begin{tabular}{C{2.5cm} C{2.5cm} L{8cm}}
\toprule
\textbf{Champ (16\,bits)} & \textbf{Valeur} & \textbf{Signification} \\
\midrule
\code{flags} & \hex{001B} & Voir tableau ci-dessous \\
\code{pins}  & \hex{000F} & Mapping broches --- valeur par d\'efaut NEO-6M \\
\bottomrule
\end{tabular}
\end{center}

\begin{infobox}[title={D\'ecomposition de flags = 0x001B}]
\begin{tabular}{C{1.5cm} C{2.5cm} L{9cm}}
\toprule
\textbf{Bit} & \textbf{Nom} & \textbf{Fonction} \\
\midrule
0 & \code{svcs}      & Alimentation DC + commutateur RF vers antenne externe \\
1 & \code{scd}       & D\'etection de court-circuit \\
3 & \code{pdwnOnSCD} & Coupe l'alimentation si court-circuit d\'etect\'e \\
4 & \code{recovery}  & R\'etablissement automatique apr\`es court-circuit \\
\midrule
\multicolumn{3}{c}{\hex{01}\,|\,\hex{02}\,|\,\hex{08}\,|\,\hex{10} = \hex{1B}} \\
\bottomrule
\end{tabular}
\end{infobox}

\subsubsection{Code source}

\begin{cppbox}
static const uint8_t kCfgAnt[] = {
    0xB5, 0x62,   // Sync UBX
    0x06, 0x13,   // Classe CFG (0x06), ID ANT (0x13)
    0x04, 0x00,   // Longueur payload : 4 octets
    0x1B, 0x00,   // flags (LE) : svcs + scd + pdwnOnSCD + recovery
    0x0F, 0x00,   // pins  (LE) : mapping broches par defaut NEO-6M
    0x47, 0x57    // Checksum Fletcher
};
\end{cppbox}

% ════════════════════════════════════════════════════════════════════════════
\section{Explication ligne par ligne de \texttt{GpsUart::begin()}}
% ════════════════════════════════════════════════════════════════════════════

Voici le code complet de la fonction, suivi d'une explication d\'etaill\'ee
de chaque bloc.

\begin{cppbox}
void GpsUart::begin() {
    gsvTotal_.begin(gps_, "GPGSV", 3);
    Serial1.begin(BoardConfig::GPS_BAUD_RATE, SERIAL_8N1,
                  BoardConfig::GPS_RX_PIN, BoardConfig::GPS_TX_PIN);
    DBG("GPS", "Serial1 started %u baud RX=GPIO%d TX=GPIO%d",
        (unsigned)BoardConfig::GPS_BAUD_RATE,
        BoardConfig::GPS_RX_PIN, BoardConfig::GPS_TX_PIN);

    delay(500);   // attendre fin du boot interne du NEO-6M

    static const uint8_t kCfgCfgReset[] = {
        0xB5, 0x62, 0x06, 0x09, 0x0D, 0x00,
        0x1F, 0x00, 0x00, 0x00,
        0x00, 0x00, 0x00, 0x00,
        0x1F, 0x00, 0x00, 0x00,
        0x17,
        0x71, 0xFE
    };
    Serial1.write(kCfgCfgReset, sizeof(kCfgCfgReset));
    DBG("GPS", "CFG-CFG sent -- factory reset to ROM defaults");
    delay(1500);  // attendre redemarrage apres reinitialisation

    while (Serial1.available()) Serial1.read();  // vider le buffer

    static const uint8_t kCfgAnt[] = {
        0xB5, 0x62, 0x06, 0x13, 0x04, 0x00,
        0x1B, 0x00,
        0x0F, 0x00,
        0x47, 0x57
    };
    Serial1.write(kCfgAnt, sizeof(kCfgAnt));
    delay(100);
    DBG("GPS", "CFG-ANT sent -- active antenna supervisor + bias enabled");
}
\end{cppbox}

\subsection{Bloc 1 --- Initialisation du parser NMEA personnalis\'e}

\begin{cppbox}
gsvTotal_.begin(gps_, "GPGSV", 3);
\end{cppbox}

\code{TinyGPSPlus} ne d\'ecode pas toutes les phrases NMEA automatiquement.
La phrase \code{GPGSV} (satellites en vue) n'est pas trait\'ee par d\'efaut~:
il faut d\'eclarer explicitement quel champ extraire.

\code{gsvTotal\_} est un objet \code{TinyGPSCustom}. Cette ligne lui dit~:
dans les phrases \code{GPGSV}, extraire le \textbf{champ num\'ero 3} (le
nombre total de satellites en vue). Ce champ devient lisible via
\code{gsvTotal\_.value()}.

\subsection{Bloc 2 --- Ouverture du port s\'erie}

\begin{cppbox}
Serial1.begin(BoardConfig::GPS_BAUD_RATE, SERIAL_8N1,
              BoardConfig::GPS_RX_PIN, BoardConfig::GPS_TX_PIN);
\end{cppbox}

Configure le port UART1 de l'ESP32 pour communiquer avec le NEO-6M~:

\begin{itemize}
  \item \code{GPS\_BAUD\_RATE} = 9\,600 bauds --- vitesse par d\'efaut ROM du NEO-6M.
  \item \code{SERIAL\_8N1} --- 8 bits de donn\'ees, pas de parit\'e, 1 bit de stop.
  \item Les broches GPIO sont d\'efinies dans \code{BoardConfig.h} et
        correspondent au c\^ablage du T-Beam.
\end{itemize}

\subsection{Bloc 3 --- Attente de d\'emarrage du GPS}

\begin{cppbox}
delay(500);
\end{cppbox}

Le NEO-6M a besoin d'environ 100 \`a 400\,ms pour terminer sa s\'equence
de mise sous tension. Sans ce d\'elai, les commandes UBX envoy\'ees
imm\'ediatement pourraient \^etre ignor\'ees. On attend 500\,ms pour une
marge confortable.

\subsection{Bloc 4 --- Envoi de CFG-CFG}

\begin{cppbox}
Serial1.write(kCfgCfgReset, sizeof(kCfgCfgReset));
\end{cppbox}

Envoie les 21 octets de la trame CFG-CFG sur le UART. \code{sizeof()} est
\'evalu\'e par le compilateur~: sa valeur est exacte et ne peut pas \^etre
oubli\'ee ou erron\'ee.

Le NEO-6M re\c{c}oit la trame, v\'erifie le checksum, puis~:
\begin{enumerate}
  \item Efface les sections sp\'ecifi\'ees en flash et BBR.
  \item Recharge ces sections depuis sa ROM (param\`etres d'usine).
  \item Envoie une trame ACK-ACK de confirmation.
  \item Red\'emarre son moteur de navigation interne.
\end{enumerate}

\subsection{Bloc 5 --- Attente du red\'emarrage}

\begin{cppbox}
delay(1500);
\end{cppbox}

La r\'einitialisation ROM prend entre 500\,ms et 1\,200\,ms. On attend
1\,500\,ms pour s'assurer que le module est de nouveau pr\^et avant d'envoyer
CFG-ANT.

\subsection{Bloc 6 --- Vidage du buffer}

\begin{cppbox}
while (Serial1.available()) Serial1.read();
\end{cppbox}

Pendant le d\'elai de 1\,500\,ms, le GPS a pu envoyer des donn\'ees sur le
UART (ACK-ACK, premi\`eres phrases NMEA de d\'emarrage). Ces donn\'ees
s'accumulent dans le buffer de l'ESP32. On les lit et on les jette pour
repartir sur une base propre avant CFG-ANT.

\subsection{Bloc 7 --- Envoi de CFG-ANT}

\begin{cppbox}
Serial1.write(kCfgAnt, sizeof(kCfgAnt));
delay(100);
\end{cppbox}

Envoie les 12 octets de la trame CFG-ANT. Le NEO-6M l'applique
imm\'ediatement~:

\begin{enumerate}
  \item Alimentation DC activ\'ee sur u.FL $\Rightarrow$ LNA Taoglas aliment\'e.
  \item Commutateur RF du T-Beam bascul\'e vers l'antenne externe.
  \item D\'etection de court-circuit activ\'ee.
\end{enumerate}

Le \code{delay(100)} laisse le temps au module d'appliquer la configuration
et de renvoyer l'ACK avant la fin de la fonction.

% ════════════════════════════════════════════════════════════════════════════
\section{Fonctionnement de \texttt{GpsUart::update()}}
% ════════════════════════════════════════════════════════════════════════════

La fonction \code{update()} est appel\'ee \`a chaque cycle de la boucle
principale (environ 50 fois par seconde). Elle lit les octets disponibles
sur le UART et les passe au parser \code{TinyGPSPlus}.

\begin{cppbox}
void GpsUart::update() {
    while (Serial1.available()) {
        char c = Serial1.read();
        gps_.encode(c);               // (1) passe l'octet au parser

        if (c == '\n') {              // (2) fin de phrase NMEA
            lineBuf_[lineLen_] = '\0';
            memcpy(completedLine_, lineBuf_, lineLen_ + 1);
            lineLen_ = 0;
        } else if (c != '\r' && lineLen_ < sizeof(lineBuf_) - 1) {
            lineBuf_[lineLen_++] = c; // (3) accumule les caracteres
        }
    }

    const bool nowValid = gps_.location.isValid(); // (4) fix GPS actif ?

    if (nowValid && !prevValid_) { DBG("GPS", "FIX ACQUIRED ..."); }
    if (!nowValid && prevValid_) { DBG("GPS", "FIX LOST"); }
    prevValid_ = nowValid;

    pos_.valid      = nowValid;       // (5) mise a jour structure publique
    pos_.lat        = nowValid ? gps_.location.lat()  : 0.0;
    pos_.lon        = nowValid ? gps_.location.lng()  : 0.0;
    pos_.speedKmph  = gps_.speed.isValid()
                        ? (float)gps_.speed.kmph()  : 0.0f;
    pos_.courseDeg  = gps_.course.isValid()
                        ? (float)gps_.course.deg()  : 0.0f;
    pos_.satellites = gps_.satellites.isValid()
                        ? (uint8_t)gps_.satellites.value() : 0;
    pos_.hdop       = gps_.hdop.isValid()
                        ? (float)gps_.hdop.hdop()   : 99.9f;
    pos_.ageMs      = gps_.location.isValid()
                        ? gps_.location.age() : 0xFFFFFFFFUL;
}
\end{cppbox}

\begin{enumerate}
  \item \textbf{\code{gps\_.encode(c)}} --- Le c\oe{}ur du parser. \code{TinyGPSPlus}
        re\c{c}oit un octet \`a la fois et reconstruit les phrases NMEA en interne.
        Quand une phrase compl\`ete et valide est re\c{c}ue, il met \`a jour ses
        champs internes (\code{location}, \code{speed}, \code{satellites}, etc.).

  \item \textbf{Capture de la phrase compl\`ete} --- Quand le caract\`ere
        \code{\textbackslash{}n} (fin de ligne NMEA) est re\c{c}u, on copie le
        buffer dans \code{completedLine\_} pour permettre l'affichage de la
        derni\`ere phrase NMEA lors du diagnostic.

  \item \textbf{Accumulation caract\`ere par caract\`ere} --- Les caract\`eres
        sont accumul\'es dans \code{lineBuf\_}. Le \code{\textbackslash{}r}
        (retour chariot) est ignor\'e~: le protocole NMEA utilise \code{CR+LF}
        comme fin de ligne, mais seul \code{LF} d\'elimite la phrase.

  \item \textbf{\code{gps\_.location.isValid()}} --- Retourne \code{true}
        uniquement si une phrase NMEA valide avec fix actif a \'et\'e re\c{c}ue
        r\'ecemment. Un fix froid (premier d\'emarrage) prend typiquement
        30 secondes \`a 2 minutes \`a l'ext\'erieur avec une bonne vue du ciel.

  \item \textbf{Mise \`a jour de \code{pos\_}} --- La structure publique est
        rafra\^{\i}chie \`a chaque appel. Un champ invalide retourne une valeur
        de s\'ecurit\'e~: \code{0.0} pour les coordonn\'ees, \code{99.9} pour
        HDOP (valeur maximale = pas de fix), \code{0xFFFFFFFF} pour l'\^age
        (signifiant donn\'ee absente).
\end{enumerate}

% ════════════════════════════════════════════════════════════════════════════
\section{Guide de diagnostic}
% ════════════════════════════════════════════════════════════════════════════

\begin{center}
\begin{longtable}{L{3.5cm} L{5.5cm} L{5.5cm}}
\toprule
\textbf{Sympt\^ome} & \textbf{Cause probable} & \textbf{Action} \\
\midrule
\endhead
\code{chars} fig\'e \`a 0
  & UART mal c\^abl\'e, broches TX/RX \'echang\'ees, ou GPS non aliment\'e
  & V\'erifier c\^ablage et alimentation \\[0.5em]
\code{chars} fig\'e \`a 10
  & Config flash corrompue --- NMEA d\'esactiv\'e
  & D\'ej\`a corrig\'e dans \code{begin()} par CFG-CFG \\[0.5em]
\code{chars} monte,\newline \code{visible} = 0
  & Obstruction physique (toit, b\^atiment)
  & Aller \`a l'ext\'erieur, ciel d\'egag\'e \\[0.5em]
\code{chars} monte,\newline \code{visible} > 0, pas de fix
  & Temps d'acquisition froid normal
  & Attendre 1 \`a 3\,min \`a l'ext\'erieur \\[0.5em]
\code{chars} monte,\newline \code{visible} = 0 m\^eme dehors
  & LNA non aliment\'e ou mauvais connecteur
  & V\'erifier CFG-ANT, c\^able u.FL et sertissage \\[0.5em]
Fix acquis puis perdu r\'eguli\`erement
  & HDOP \'elev\'e, obstruction partielle
  & D\'egager la vue du ciel (m\^at, pont) \\
\bottomrule
\end{longtable}
\end{center}

% ════════════════════════════════════════════════════════════════════════════
\section{Chronologie de l'initialisation}
% ════════════════════════════════════════════════════════════════════════════

\begin{center}
\begin{tabular}{C{2.5cm} L{11cm}}
\toprule
\textbf{Temps (ms)} & \textbf{\'Ev\'enement} \\
\midrule
0             & \code{Serial1.begin()} --- UART ESP32 ouvert \`a 9\,600 bauds \\
0 -- 500      & D\'elai de d\'emarrage --- le NEO-6M termine son boot interne \\
500           & Envoi de CFG-CFG --- demande de r\'einitialisation ROM \\
$\approx$ 501 & Le NEO-6M r\'epond ACK-ACK (10 octets) \\
501 -- 2000   & D\'elai de red\'emarrage --- rechargement de la config ROM \\
2000          & Vidage du buffer UART \\
2000          & Envoi de CFG-ANT --- activation du superviseur d'antenne \\
$\approx$ 2100 & Fin de \code{begin()} --- le GPS est pr\^et \`a \'emettre du NMEA \\
2100+         & \code{update()} appel\'e en boucle --- d\'ecodage des phrases NMEA \\
$\approx$ 60\,000 -- 180\,000
              & Premier fix GPS (typique en ext\'erieur, ciel d\'egag\'e) \\
\bottomrule
\end{tabular}
\end{center}

% ════════════════════════════════════════════════════════════════════════════
\section*{Conclusion}
\addcontentsline{toc}{section}{Conclusion}
% ════════════════════════════════════════════════════════════════════════════

Le dysfonctionnement du GPS \'etait d\^u \`a la persistance d'une configuration
corrompue dans la m\'emoire flash interne du NEO-6M, qui avait d\'esactiv\'e
toutes les sorties NMEA. La correction consiste \`a envoyer syst\'ematiquement,
\`a chaque d\'emarrage, une commande de r\'einitialisation usine (\textbf{CFG-CFG})
suivie de l'activation du superviseur d'antenne active (\textbf{CFG-ANT}).

Ces deux commandes sont int\'egr\'ees dans \code{GpsUart::begin()}, ce qui
garantit un comportement correct ind\'ependamment de toute session de
configuration ant\'erieure effectu\'ee avec u-center ou tout autre outil.

\begin{okbox}[title=\'Etat actuel du syst\`eme]
\begin{itemize}
  \item[$\checkmark$] Phrases NMEA r\'eactiv\'ees syst\'ematiquement \`a chaque boot.
  \item[$\checkmark$] Antenne active Taoglas ADFGP.25A aliment\'ee correctement.
  \item[$\checkmark$] Commutateur RF T-Beam bascul\'e vers l'antenne externe.
  \item[$\checkmark$] Fix GPS confirm\'e lors des tests en ext\'erieur.
\end{itemize}
\end{okbox}

\end{document}
